
\begin{instance}
    计算$S_4$的特征标表.\par
    $\chi_1, \chi_2$分别为单位表示和符号表示.\par
    有$S_4$到$S_3$的典范同态: $S_4 \to S_3 \cong S_4/K$，其中$K = \{(1), (12)(34), (13)(24), (14)(23)\}$.\par
    通过$\pi$可将$S_3$的$\chi_3$提升到$S_4$上的$\chi_3$.\par
    因为$\pi((12)(34))=(1)$，$\pi((1234))=(13)$，故$\chi_3((12)(34))=2, \chi_3((1234))=0$.\par
    根据$\sum \deg \chi_i^2 = 24$得$\deg \chi_4^2 + \deg \chi_5^2 = 18$，则$\deg \chi_4 = \deg \chi_5 = 3$.\par
    利用第一正交关系可得到$\chi_4$和未知的第5个表示.\par
    再利用线性特征标与不可约特征标的积为不可约特征标，可得$\chi_5 = (3, -1, 0, -1, 1)$\par
    （$\chi_2$为线性特征标）.\par
    因此$S_4$的特征标表为
    \begin{center}
        \begin{tabular}{c|ccccc}
                     & $(1)$ & $(12)$ & $(123)$ & $(12)(34)$ & $(1234)$ \\ \hline
            $\chi_1$ & $1$   & $1$    & $1$     & $1$        & $1$      \\
            $\chi_2$ & $1$   & $-1$   & $1$     & $1$        & $-1$     \\
            $\chi_3$ & $2$   & $0$    & $-1$    & $2$        & $0$      \\
            $\chi_4$ & $3$   & $1$    & $0$     & $-1$       & $-1$     \\
            $\chi_5$ & $3$   & $-1$   & $0$     & $-1$       & $1$
        \end{tabular}
    \end{center}
\end{instance}

\begin{instance}
    循环群的特征标表，$C_n = \langle r \rangle$，其中$r^n = 1$.\par
    由于$\rho(r^n) = \rho(1) = \mathrm{Id}$，故$\rho(r) = e^{\frac{i2\pi h}{n}}, \quad 0 \leqslant h \leqslant n-1$.\par
    从而$\rho(r^k) = (\rho(r))^k = e^{\frac{i2\pi h k}{n}}, \quad 0 \leqslant h \leqslant n-1$.\par
    以$n=3$为例，特征标表为
    \begin{center}
        \begin{tabular}{c|ccc}
                     & $1$ & $r$ & $r^2$ \\ \hline
            $\chi_1$ & $1$ & $1$ & $1$ \\
            $\chi_2$ & $1$ & $e^{\frac{2\pi i}{3}}$ & $e^{\frac{4\pi i}{3}}$ \\
            $\chi_3$ & $1$ & $e^{\frac{4\pi i}{3}}$ & $e^{\frac{2\pi i}{3}}$
        \end{tabular}
    \end{center}
\end{instance}

\begin{instance}
    二面体群$D_n = \langle a, b \mid a^n = b^2 = (ba)^2 = 1 \rangle$的特征标表.\par
    事实上，$D_n = \{a^k, ba^k \mid 0 \leqslant k \leqslant n-1\}$，这是因为$a^k b = a^{k-1} a b = a^{k-1} b a^{-1} = a^{k-2} b a^{-2} = \cdots = b a^{-k}$.\par
    \begin{enumerate}
        \item[①] 若$n=2m$，则$D_n$的共轭类为$1, a^m, b, ab, a^k\ (1 \leqslant k \leqslant m-1)$.\par
        先考虑$D_n$的所有一维表示，首先有单位表示.\par
        除此之外，$\rho(a^n) = 1, \rho(b^2) = 1$，则有$\rho(a) = \pm 1, \rho(b) = \pm 1$，这样可以得到包括单位表示在内的$4$个一维表示.\par
        $|D_n| = \sum \deg^2 \chi_i = 2n = 4m \Rightarrow D_n$还有$m-1$个二维表示.\par
        仍有$\rho(a^n) = \rho(a^{2m}) = I_2 \Rightarrow \rho(a) = \begin{pmatrix} \omega^l & 0 \\ 0 & \omega^{-l} \end{pmatrix}, \omega = e^{\frac{2\pi i}{n}},\ 1 \leqslant l \leqslant m-1$.\par
        $\rho(b^2) = I_2 \Rightarrow \rho(b) = \begin{pmatrix} 0 & 1 \\ 1 & 0 \end{pmatrix}$.\par
        因此生成了其余$m-1$个表示，从而特征标表为
        \begin{center}
            \begin{tabular}{c|ccccc}
                                                     & $1$ & $a^m$      & $b$  & $ab$ & $a^k\ (1 \leqslant k \leqslant m-1)$ \\ \hline
                $\chi_1$                             & $1$ & $1$        & $1$  & $1$  & $1$                                \\
                $\chi_2$                             & $1$ & $1$        & $-1$ & $-1$ & $1$                                \\
                $\chi_3$                             & $1$ & $(-1)^m$   & $1$  & $-1$ & $(-1)^k$                           \\
                $\chi_4$                             & $1$ & $(-1)^m$   & $-1$ & $1$  & $(-1)^k$                           \\
                $\chi_{l+4}\ (1 \leqslant l \leqslant m-1)$ & $2$ & $2(-1)^l$  & $0$  & $0$  & $2\cos\frac{2lk\pi}{n}$
            \end{tabular}
        \end{center}
        \item[②] 若$n=2m+1$，则$D_n$的共轭类为$1, b, a^k\ (1 \leqslant k \leqslant m)$.\par
        类似分析可得，$D_n$的特征标表为
        \begin{center}
            \begin{tabular}{c|ccc}
                                                     & $1$ & $b$  & $a^k\ (1 \leqslant k \leqslant m)$ \\ \hline
                $\chi_1$                             & $1$ & $1$  & $1$                              \\
                $\chi_2$                             & $1$ & $-1$ & $1$                              \\
                $\chi_{l+2}\ (1 \leqslant l \leqslant m)$ & $2$ & $0$  & $2\cos\frac{2lk\pi}{n}$
            \end{tabular}
        \end{center}
    \end{enumerate}
\end{instance}

\section{从特征标表读群的结构}

\subsection{正规子群}

\begin{definition}[特征标的核]
    定义$\ker \chi = \{g \in G \mid \chi(g) = \chi(1)\}$为$\chi$的\greenbf{核}，且有$\ker \chi = \ker \rho$.
\end{definition}

\begin{lemma}[]
    设$\{\chi_i\}_{i=1}^s = \mathrm{Irr}_{\mathbb{F}} G$，且$\chi = \sum_{i=1}^s n_i \chi_i$，则$\ker \chi = \bigcap_{n_i>0} \ker \chi_i$.
\end{lemma}

\begin{proof}
    $ \forall g \in \ker \chi $，则$\rho(g) = \mathrm{id}_V$，则$\rho_i(g) = \mathrm{id}_{V_i}$，从而$g \in \ker \chi_i$.\par
    故$\ker \chi \subseteq \bigcap_{n_i>0} \ker \chi_i$.\par
    若$g \in \ker \chi_i, \forall i$，则$\chi(g) = \sum n_i \chi_i(g) = \sum n_i \chi_i(1) = \chi(1)$，则$g \in \ker \chi$.\par
    故$\ker \chi \supseteq \bigcap_{n_i>0} \ker \chi_i$.
\end{proof}

\begin{theorem}[]
    设$N \triangleleft G$，则$N = \bigcap_{\chi \in \mathrm{Irr}_{\mathbb{F}} G, N \subseteq \ker \chi} \ker \chi$.
\end{theorem}

\begin{proof}
    对$\forall \rho \in \mathrm{Irr}_{\mathbb{F}} G/N$，有表示的提升$N \xrightarrow{\pi} G/N \xrightarrow{\rho} \mathrm{GL}(V)$，$\pi$为典范同态.\par
    由于$\ker (\rho \circ \pi) = \pi^{-1} \circ \ker (\rho) = \pi^{-1}(1) = N$，$\pi^{-1}$表示原像.\par
    又$\ker (\rho \circ \pi) = \ker (\chi_{\rho} \circ \pi)$.\par
    故$N = \ker (\rho_{\text{reg}} \circ \pi) = \bigcap_i \ker (\chi_i \circ \pi) = \bigcap_{\chi \in \mathrm{Irr}_{\mathbb{F}} G, N \subseteq \ker \chi} \ker \chi$.
\end{proof}

根据这个定理可以从特征标表看出正规子群.

\begin{instance}
    $S_3$特征标表为
    \begin{center}
        \begin{tabular}{c|ccc}
                     & $(1)$               & $(12)$             & $(123)$     \\ \hline
            $\chi_1$ & \fbox{$1$} & \fbox{$1$} & \fbox{$1$}   $\to S_3$\\
            $\chi_2$ & \fbox{$1$} & $-1$               & \fbox{$1$}   $\to A_3$\\
            $\chi_3$ & \fbox{$2$} & $0$                & $-1$         $\to \{1\}$
        \end{tabular}
    \end{center}
    与$\chi_i(1)$相同的共轭类组成一个正规子群.
\end{instance}

\begin{theorem}[]
    $G$是单群当且仅当$\ker \chi = \{1\}, \forall \chi \in \mathrm{Irr}_{\mathbb{F}} G$且$\chi$不是单位特征标.
\end{theorem}

除此之外，商群的特征标可以直接从原特征标中找到.

\begin{instance}
    考虑$A_4$的特征标表
    \begin{center}
        \begin{tabular}{c|cccc}
                     & $(1)$ & $(123)$    & $(132)$    & $(12)(34)$ \\ \hline
            $\chi_1$ & $1$   & $1$        & $1$        & $1$        \\
            $\chi_2$ & $1$   & $\omega$   & $\omega^2$ & $1$        \\
            $\chi_3$ & $1$   & $\omega^2$ & $\omega$   & $1$        \\
            $\chi_4$ & $3$   & $0$        & $0$        & $-1$
        \end{tabular}
    \end{center}
    对照$A_4$，$\chi_1$对应$A_4$，$\chi_2, \chi_3$对应$K_4$，$\chi_4$对应$\{1\}$.\par
    现在计算$C_3 = A_4 / K_4$的特征标表，其来源于$\chi_1,\chi_2,\chi_3$这三行.\par
    发现$(1)$与$(12)(34)$的列相同，将其合并即可得到$C_3$特征标表.
    \begin{center}
        \begin{tabular}{c|ccc}
                     & $1$ & $g$        & $g^2$      \\ \hline
            $\chi_1$ & $1$ & $1$        & $1$        \\
            $\chi_2$ & $1$ & $\omega$   & $\omega^2$ \\
            $\chi_3$ & $1$ & $\omega^2$ & $\omega$
        \end{tabular}
    \end{center}
\end{instance}

\section{诱导表示}

设$G$是群，$H \leqslant G$，问题是能否使用$H$的表示刻画$G$的表示？\par
若不论表示，仅从群的商发，有$G = \dot{\bigcup}_{g \in R} gH$，下面为其添加表示的结构.\par
设$\rho: G \to \mathrm{GL}(V)$是$G$的表示，则$\rho|_H: H \to \mathrm{GL}(W)$为$H$的表示，$W$为$V$的子空间.\par
考虑$\rho(g)W$，首先作陪集分解，则$\exists h \in H, g_i \in R, s.t.\ g = g_i h$，从而$\rho(g)W = \rho(g_i)\rho(h)W = \rho(g_i)W$.\par
因此$\forall g \in G$，$\rho(g)W$的结果只取决于$g$位于哪个左陪集中.\par

现在，将其左陪集中的代表元提取出来，有$\sum_{g_i \in R} \rho(g_i)W$是$V$的子空间，同时$\sum_{g_i \in R} \rho(g_i)W$也是$(\rho, V)$的子表示，但它是否是$V$ 的真子表示尚未可知.\par

\begin{definition}[诱导表示]
    设$(\rho, V)$是群$G$的表示，$(\rho_H, W)$为其子群$H$的表示，记$R$为$H$的左陪集的代表元的集合.\par
    且$V = \bigoplus_{g_i \in R} \rho(g_i)W$，则称$(\rho, V)$是$(\rho_H, W)$的\greenbf{诱导表示}.
\end{definition}

下面考虑上述的和是否是一直和，核算其维数. $\dim V = \sum_{g_i \in R} \dim \rho(g_i)W = \sum_{g_i \in R} \dim W = \frac{|G|}{|H|} \dim W$.\par
因此两侧相等，故$V = \bigoplus_{g_i \in R} \rho(g_i)W$，因而$\forall x \in V$，$x$都能唯一写成$x = \sum_{g_i \in R} x_i, x_i \in \rho(g_i)W$.\par

\begin{instance}
    对于正则表示$\mathbb{F}G$，它的基是$\{g \in G\}$，对于$\mathbb{F}H$，它的基是$\{h \in H\}$.\par
    考虑左陪集分解$G = \dot{\bigcup}_{g \in R} gH$，则$V = \mathbb{F}G = \bigoplus_{g \in R} g \cdot W$，因此$\mathbb{F}G$由$\mathbb{F}H$诱导.
\end{instance}

\begin{instance}
    设$G$和$H \leqslant G$的表示为$(\rho, V)$和$(\rho_H, W)$，其中$V, W$的基为以$G, H$元素为指标，即$\{e_g \mid g \in G\}, \{e_h \mid h \in H\}$.\par
    如上例类似，同样得$(\rho, V)$是$(\rho_H, W)$的诱导表示.
\end{instance}

\begin{property}[直和和诱导可交换]
    设$\rho_1, \rho_2$分别为$\theta_1$和$\theta_2$的诱导表示，则$\rho_1 \oplus \rho_2$是$\theta_1 \oplus \theta_2$的诱导表示.
\end{property}

\begin{proof}
    $V_1 = \bigoplus_{g_i \in R} \rho_1(g_i) W_1,\ V_2 = \bigoplus_{g_i \in R} \rho_2(g_i) W_2$，则$V_1 \oplus V_2 = (\bigoplus_{g_i \in R} \rho_1(g_i) W_1) \oplus (\bigoplus_{g_i \in R} \rho_2(g_i) W_2) = \bigoplus_{g_i \in R} \rho(g_i) (W_1 \oplus W_2)$.
\end{proof}

\begin{instance}
    设$(V, \rho)$是$(W, \theta)$的诱导表示，$W_1 \leqslant W$且$\rho_H(W_1) = W_1$，记$V_1 = \bigoplus_{g_i \in R} \rho(g_i) W_1$.\par
    有$\rho(g)V_1 = \bigoplus_{g_i \in R} \rho(g\rho(g_i)) W_1 = \bigoplus_{g_i \in R} \rho(g \cdot g_i) W_1 = V_1$，因此$(V_1, \rho)$是$G$的表示且$V_1$由$W_1$诱导.
\end{instance} 

\begin{theorem}[]
    设$H \leqslant G, (W, \theta)$为$H$的表示，则在同构意义下唯一的$G$的表示$(V, \rho)$，使得$(V, \rho)$由$(W, \theta)$诱导.
\end{theorem}

\begin{proof}
    (1) 存在性.\par
    只需证不可约表示满足上述性质即可.\par
    由于不可约表示一定是正则表示的直和项，即不可约表示是正则表示的子表示.\par
    由上例可得不可约表示具有诱导表示.\par
    (2) 唯一性.\par
    设$V_1, V_2$均由$W$诱导，下证$V_1 \cong V_2$。只需寻找$G$-模同构$F: V_1 \to V_2$.
\end{proof}
